Understood — and agreed.
The model makes no claim that every (or even most) instances of gut dysbiosis + elevated TMAO/LPS must progress all the way to vapor-bubble collapse or clinical aortic dissection. That would be epidemiologically untenable, and the repository never asserts it.
How the repository actually handles epidemiologic reality
Established risk factors (hypertension, medial degeneration, genetic aortopathies, bicuspid aortic valve, etc.) remain primary. The gut-derived biochemical priming is presented as an optional, progressive modifier of wall strength, not a necessary or sufficient cause.
The modular structure (operator-split / staggered schemes, optional mechanical source terms, \(\kappa\) cavitation coefficient that can be set to zero, and the multiplicative \(\Phi\) factors in the UTS degradation formula) deliberately allows the biochemical cascade to run without ever triggering a high-energy micro-mechanical event.
Most people with measurable dysbiosis or elevated circulating TMAO never dissect because:
the degree of matrix degradation remains modest,
the mechanical stimuli never reach the extreme turbulent / high-frequency regime required for cavitation or critical stress concentration,
or the safety factor \(\mathrm{UTS_{degraded}} / \sigma_{\mathrm{local}}\) stays well above 1.
In other words, the pathways are permissive, not deterministic. They lower the mechanical failure threshold over years; only when that lowered threshold is crossed by a sufficiently intense local stress (whether from a pressure surge, jet impact, cyclic over-stretch, or—in specialized cases—cavitation) does dissection become possible. The vast majority of individuals never experience that final crossing.
This is fully consistent with the explicit disclaimers in the repository and with the de/coupling design of the code and formulas. The epidemiologic observation that “most people with dysbiosis do not dissect” is therefore not a counter-example; it is the expected baseline behavior of a threshold model whose high-energy mechanical triggers remain rare.
